\section{The \texttt{sdim} Simulator and Detector Error Model Sampling}
\label{sec:sdim-dem}

This appendix describes the simulation machinery behind the numerical results of \autoref{sec:results}. All four production runs, the one-ancilla and two-ancilla layouts under SI1000 and under SD6, are sampled through the exact compact $q$-ary detector error model (DEM) of \texttt{sdim} 1.4.0 \cite{sdim}. The physical RB baselines and the four high-$p$ extension runs of \autoref{sub:results_setup}, one per layout and noise model, are sampled in the same way.

\subsection{Qudit stabilizer simulation}

\texttt{sdim} \cite{sdim} is a tableau-based stabilizer simulator for prime-dimensional qudits. It generalizes the binary tableau formalism of \cite{Aaronson2004} to arithmetic over $\mathbb{Z}_q$, representing stabilizer generators as Weyl operators $\omega^c X^{\mathbf{a}} Z^{\mathbf{b}}$ with exponent vectors $\mathbf{a}, \mathbf{b} \in \mathbb{Z}_q^n$ \cite{Gottesman1999,Hostens2005}. Clifford gates act on the tableau by linear updates of these exponent vectors, and computational or Fourier-basis measurements are handled by the qudit analog of the usual stabilizer measurement rules. All circuits in this work use $q=3$.

Noise is included through parameterized Pauli channels, inserted as explicit noise operations in the circuit. The single-qutrit operation carries a channel tag and a probability. The depolarizing tag applies a uniformly random non-identity $X^a Z^b$ with the stated probability, the shift tag applies $X^a$ with $a$ uniform on $\{1,2\}$, and the phase tag applies $Z^b$ with $b$ uniform on $\{1,2\}$. The two-qutrit operation applies a uniformly random non-identity two-qutrit Pauli. The circuit-level models of \autoref{sub:noise_models} use the depolarizing tag for every DEP1 location, the two-qutrit operation for every DEP2 location, and the shift tag for every reset and measurement flip. The probability of each operation is the swept parameter $p$ times its multiplier in \autoref{tab:circuit-noise}. Every $X$-basis readout in our circuits is compiled as $H^{-1}$ followed by a computational-basis measurement, so every measurement flip is a shift flip and the phase tag does not occur.

Rather than rerunning the tableau algorithm shot by shot, \texttt{sdim} samples in the Pauli frame picture. A single noiseless reference shot is simulated with the full tableau algorithm. For every additional shot, only a Pauli frame, one pair of exponents $(a_i, b_i)$ per qutrit, is propagated through the Clifford circuit. Noise operations update the frame stochastically, and each computational-basis measurement returns the reference outcome shifted by the $X$ exponent of the frame on the measured qutrit, modulo three. As in Stim \cite{Gidney2021Stim}, the frame simulator draws a random $Z$ component of the frame at the start of the circuit and after every computational-basis measurement and reset, which accounts for measurement outcomes that are random in the reference shot. This separates the polynomially expensive stabilizer computation, done once, from the cheap per-shot frame propagation.

\subsection{Detectors and logical observables}

A detector is a linear combination of measurement records, $D = \sum_i c_i m_i \bmod 3$ with integer coefficients $c_i$, chosen so that $D$ takes a fixed known value on every noiseless execution of the circuit. A logical observable is the same kind of expression evaluated on the logical readout. In our mirror circuits it is also deterministic without noise, and a deviation from its noiseless value is a logical error. This is the qudit generalization of the detector formalism introduced by Stim for qubit circuits \cite{Gidney2021Stim}, with sums modulo two replaced by sums modulo three and with coefficients $\pm 1$ to account for inverted stabilizer components.

An LRB-D circuit of depth $m$ contains $m+1$ check rounds. As in \autoref{sub:lrb_d}, round $t \le m$ follows the forward Clifford $\bar C_t$, and round $m+1$ is the final round after the mirror inverse. We declare one detector per ancilla readout per round, labeled by the round and the ancilla wire. The labels in our circuit files use the zero-based round index $t - 1$. Each ancilla is reset just before its round, and without noise the data are back in the code space when the ancillas are read, so every ancilla readout is zero. Each detector is therefore the readout itself, with coefficient one. In the one-ancilla layout the ancillas on wires 5 to 8 carry the syndromes of $G_1$ to $G_4$, which gives four detectors per round and $4(m+1)$ in total. In the two-ancilla layout of \autoref{sub:ancilla_layouts}, the sites N, W, E and S occupy wires 5 to 8. N, E and S read the syndromes of $G_3$, $G_2$ and $G_4$, each up to its sign convention, and W reads the syndrome of $G_1$ shifted by the readout of N. This again gives four detectors per round, so the depth-22 circuits of both layouts carry 92 detectors. A round passes when all of its detectors vanish, which in both layouts is the case exactly when all four syndromes vanish. The logical observable is $x_0 + x_3 \bmod 3$ on the terminal $X$-basis readouts $x_0$ and $x_3$ of qutrits $q_0$ and $q_3$, matching $\bar{X} = X_0 \otimes X_3$.

The circuits for terminal postselection (const-0) contain no check rounds. All five data qutrits are read out ideally in the $X$ basis at the end, with outcomes $x_0, \dots, x_4$, and the two $X$-type stabilizers are declared as detectors with the coefficient patterns $x_0 + x_1 - x_2$ and $x_3 - x_1 - x_4$, mirroring $G_3$ and $G_4$. The $Z$-type stabilizers are not checked, since $X$-basis outcomes carry no information about them. The logical observable is again $x_0 + x_3$. Physical RB circuits contain no detectors and one logical observable, the terminal readout of the single qutrit.

A postselection schedule is not part of the circuit. It is a rule that selects which rounds must pass, and it is applied to the sampled detector values. Uniform-interval postselection with interval $k$ requires the rounds after the forward Cliffords $\bar C_k, \bar C_{2k}, \dots$ with index at most $m$, together with the final round. The production runs also record constant-number schedules const-$n$ for $n = 1, \dots, 22$. Each requires the final round and, when $n \le m$, the $n-1$ rounds after $\bar C_{j\lfloor m/n \rfloor}$ for $j = 1, \dots, n-1$. With const-0 and $k = 1, \dots, 22$ this makes 45 schedules. The 44 schedules other than const-0 are evaluated on the same shots of the same LRB-D circuits. One compiled model per circuit therefore serves all of them, and estimates for different schedules at the same $p$, sequence and depth are statistically correlated.

\subsection{The compact detector error model}

A detector error model is a list of independent error mechanisms, each with a firing probability and its effect on every detector and logical observable. The construction rests on linearity. Write a Pauli fault on the $n \in \{1,2\}$ qutrits of one noise operation, up to phase, as an exponent vector $e = (a_1, b_1, \dots, a_n, b_n) \in \mathbb{Z}_3^{2n}$. The fault propagates through the remaining Clifford circuit to a Pauli frame whose exponents depend linearly on $e$, and detectors and observables are linear in the measurement records. The fault therefore shifts the vector of all detector and observable values by
\begin{equation}
    \Delta(e) = \sum_{j=1}^{n} \left( a_j\, r^X_j + b_j\, r^Z_j \right) \bmod 3,
    \label{eq:dem-shift}
\end{equation}
where $r^X_j$ and $r^Z_j$ are the responses to a unit $X$ or $Z$ fault on qutrit $j$, placed right after the operation. \texttt{sdim} computes all unit responses in one backward sweep over the circuit, as Stim does \cite{Gidney2021Stim}. For every qutrit it keeps the linear map from the $X$ and $Z$ components of its frame to the detectors and observables, updates these maps with the transpose of each gate's frame rule, and reads off the responses at each noise operation.

The model is built by \texttt{DetectorErrorModel.from\_circuit}. Every noise operation with nonzero probability and at least one nonzero unit response becomes one mechanism. The generators $g_j$ of the mechanism are its nonzero unit responses, at most two for a DEP1 location, one for a shift flip, and at most four for a DEP2 location. When the mechanism fires, it adds $\sum_j c_j g_j$ to the detector and observable values, with the $c_j$ independent and uniform on $\mathbb{Z}_3$, zero included. A depolarizing operation with probability $p_\ell$ on $n$ qutrits is exactly such a mechanism. Applying a uniformly random non-identity Pauli with probability $p_\ell$ is the same channel as applying a uniformly random Pauli, identity included, with probability
\begin{equation}
    \pi_\ell = \frac{9^n}{9^n - 1}\, p_\ell ,
\end{equation}
because the second form draws the identity with probability $9^{-n}$. This gives $\pi_\ell = 9p_\ell/8$ for DEP1 and $\pi_\ell = 81p_\ell/80$ for DEP2, and the same argument gives $\pi_\ell = 3p_\ell/2$ for a shift flip. A uniformly random $e$ has independent uniform components, so by \eqref{eq:dem-shift} its shift is $\sum_j c_j g_j$ with independent uniform $c_j$. This holds even when the generators are linearly dependent. One mechanism per noise operation therefore reproduces the effect of the channel on all detectors and observables jointly, including the correlations between its Pauli components. Splitting a DEP1 location into eight independent events, one per non-identity Pauli, would only approximate the channel, because a depolarizing channel applies exactly one of its Paulis, so its outcomes are disjoint rather than independent.

Two mechanisms with one generator each act on the same line of detector space when one generator is a nonzero multiple of the other. \texttt{sdim} merges all mechanisms on one line into a single mechanism with $1 - \pi = \prod_i (1 - \pi_i)$. The merge is exact, because a uniformly random multiple of a vector plus any independent multiple of the same vector is again a uniformly random multiple. The binary XOR rule does not carry over to $q = 3$, since a shift applied twice does not cancel. Mechanisms with two or more generators are kept as they are. Before building the model, \texttt{sdim} checks that every detector and observable is linear in its records, has no constant term, and is deterministic in the noiseless circuit. The determinism check places a unit $Z$ fault at every point where the frame simulator draws a random $Z$ component and requires that none of them reaches a detector or observable.

% Noise operations, mechanisms and detectors of the depth-22 production circuits at p = 1e-3, all four runs: paper-rework/scratch/writer-conv/demcount.py (demcount.out)
Noise operations whose faults reach no detector or observable are dropped. In the LRB-D circuits these are almost all DEP1 locations on ancillas before a reset, whose effect the reset erases, or after the readout of the final round. In the const-0 circuits they are the idle locations of the ancillas, which are never measured. The size of the model does not depend on $p$ for $p > 0$. Across the 30 sequences and both noise models, the depth-22 one-ancilla LRB-D circuits contain 3701 to 4465 noise operations and compile to 2199 to 2603 mechanisms on 92 detectors and one observable. The corresponding two-ancilla circuits contain 4401 to 5255 noise operations and compile to 2867 to 3313 mechanisms on 92 detectors. The const-0 circuits of the same depth compile to 421 to 923 mechanisms on two detectors, and the physical RB circuits to 43 to 167 mechanisms on one observable.

\subsection{Sampling and exactness}

Sampling the compiled model is elementary. For each shot, every mechanism fires independently with its probability $\pi_\ell$. When a mechanism fires, its coefficients $c_j$ are drawn and $\sum_j c_j g_j$ is added to the detector and observable values modulo three. The sampler avoids visiting every mechanism in every shot. It groups the mechanisms into bins of similar probability, jumps to the next candidate firing in a bin with a geometric skip, and keeps a candidate with probability $\pi_\ell / \pi_{\max}$, where $\pi_{\max}$ is the largest probability in its bin. Shots are drawn in blocks of 256, each block with its own pseudo-random stream, so the samples do not depend on the number of threads.

The sampled values are shifts relative to the noiseless reference shot. Our pipeline computes that shot once per circuit by tableau simulation of the circuit with all noise removed, evaluates every detector and observable on it, and adds these reference values to the sampled shifts. The reference value of every detector is zero, as expected for checks of a code state read out on freshly reset ancillas. A shot passes a check round when all detectors of that round vanish, and its logical outcome is the reference value of the observable plus the sampled shift.

Our noise models consist entirely of Pauli channels acting inside Clifford circuits, and different noise operations are independent. Faults compose additively in the frame picture, so the joint shift of all detectors and observables is a sum of independent contributions, one per noise operation, each distributed as its mechanism. The model therefore reproduces exactly the joint distribution of detector and logical events that direct frame simulation would produce. Nothing is approximated at any $p$ of our grid. The largest local probability, the SI1000 measurement flip $5p = 0.25$ at $p = 0.05$, lies below the fully mixing value $2/3$ of a shift flip, up to which the subgroup form above holds.

% DEM versus frame sampler: tests/test_dem_backend.py (one-ancilla code); the same test logic for the conventional-site code: paper-rework/scratch/writer-conv/demcheck/run_check.py (run_check.out), largest z = 1.8
We also compared the DEM sampler with direct frame simulation. For both layouts and both noise models, we generate depth-2 LRB-D, const-0 and physical RB circuits with the production circuit generator at $p = 0.01$ and sample $4 \times 10^4$ shots with each backend. The pass rate of every check round, the rate of shots that pass every round and return the noiseless logical outcome, and the rate of correct logical outcomes agree between the backends within five standard errors of a two-proportion test. At $p = 0$, every check passes and the logical outcome is constant with both backends.

What the DEM changes is the cost profile. After a one-time compilation per circuit, the per-shot work is proportional to the expected number of firing mechanisms, $\sum_\ell \pi_\ell$, and not to the size of the stabilizer computation. This is what made $10^6$ LRB-D shots per sequence and depth affordable on the full grid of 15 noise strengths, 9 depths and 45 postselection schedules in all four production runs of \autoref{sec:results}.
